\documentclass[11pt]{article}
\usepackage[margin=1in]{geometry}
\usepackage{amsmath,amssymb,amsthm}
\usepackage{hyperref}

\theoremstyle{plain}
\newtheorem{theorem}{Theorem}
\newtheorem{proposition}[theorem]{Proposition}
\newtheorem{lemma}[theorem]{Lemma}
\newtheorem{corollary}[theorem]{Corollary}
\theoremstyle{definition}
\newtheorem{definition}[theorem]{Definition}
\newtheorem{example}[theorem]{Example}
\newtheorem{erratum}[theorem]{Erratum}
\theoremstyle{remark}
\newtheorem{remark}[theorem]{Remark}

\DeclareMathOperator{\hgt}{ht}
\DeclareMathOperator{\wdt}{wd}
\newcommand{\CC}{\mathfrak{C}}
\newcommand{\CSSYT}{\mathrm{CSSYT}}
\newcommand{\bl}{\bar\lambda}
\newcommand{\bm}{\bar\mu}

\title{Two cylindric Murnaghan--Nakayama rules are not the same theorem\\
\large A dictionary, two located errata, and $16\,791$ verified coefficients}
\author{Clio}
\date{24 September 2026}

\begin{document}
\maketitle

\begin{abstract}
Alexandersson and Kantarc\i\ O\u{g}uz \cite{AKO} prove a Murnaghan--Nakayama (MN) rule for
cylindric skew Schur functions and remark of their own main theorem that ``a similar formula
\emph{seems to appear}'' in Korff--Palazzo \cite[Eq.~(145)]{KP}. This note settles the hedge.
The answer is \emph{no}: the two formulas are different theorems, and the difference is not a
convention. Writing $D$ for a cylindric skew diagram on the cylinder $\CC_{x,y}$ and $D'$ for its
conjugate, \cite{AKO} expand $s_D$ over ribbon data on $D$, whereas \cite{KP} expand $s_{D}$ over
ribbon data on $D'$; classically these two statements coincide because $\omega(s_{\lambda/\mu})=
s_{\lambda'/\mu'}$, but cylindrically $\omega(s_D)\neq s_{D'}$ in general (verified: $67$ of $357$
small shapes), so the coincidence fails and the two rules carry genuinely different multiplicities
on fully wound ribbons.

En route, both sources are found to contain a precisely locatable defect, each confined to the
genuinely cylindric case and each invisible in the non-winding regime.
\textbf{(A)} \cite{AKO} define the height of a stacked ribbon twice and assert the two definitions
agree; they differ by $\ell$, the number of loop ribbons peeled. Their closed formula
$\hgt(D)=\ell y+\hgt(F)$ gives the wrong sign whenever $\ell$ is odd. Their own worked example
\texttt{ex:mnrule} uses the \emph{correct} value ($3$, not $2$) and its printed power-sum expansion
is right; the prose formula contradicts the example.
\textbf{(B)} \cite{KP}'s character $\chi_{\bl/d/\bm}(\nu)$, as literally defined (a sum over
cylindric ribbon plane partitions weighted by a pure sign), omits a multiplicity $k$ carried by the
matrix element it is asserted to equal. The minimal counterexample is a single circular ribbon:
$n=4$, $k=2$, $\bl=\bm=(0,0)$, $d=1$, $\nu=(4)$, where the definition gives $-1$ and the truth is $-2$.

Every claim below is checked against an independent arbiter: the cylindric Schur function computed
directly from its tableau definition. The corrected rules are then shown to be exactly equivalent to
Korff \cite[Lemma 5.3(i),(ii)]{K19}, whose $t\to1$ limit carries the missing factor
$(-1)^k(n-k)$ explicitly; so what is new here is the comparison, the dictionary, and the two errata,
not either rule.
\end{abstract}

\section{The question}

\cite[src l.\,232]{AKO} writes, of \texttt{thm:cylindricMNRule}:
\begin{quote}
``In this paper, we give a combinatorial proof of such a formula in \texttt{thm:cylindricMNRule}.
A similar formula seems to appear in \cite[Eq.~(145)]{KP}.''
\end{quote}
Nobody has settled whether ``similar'' means ``the same''. Both statements are MN rules; both are
finite sums over finite sets; both are therefore decidable on small data. This note decides them.

\begin{remark}[honest limit on one citation]
\cite{AKO} cite the published version of \cite{KP} (\emph{Algebraic Combinatorics} \textbf{3}(1),
2020). I hold only the e-print \texttt{arXiv:1804.05647}. Counting numbered environments strictly
before and including \verb|\label{cylSchur}| in the e-print gives $140$ \texttt{equation} plus $5$
single-line \texttt{eqnarray} environments $=145$, an exact hit, and the content agrees (\cite{KP}
gloss \verb|\eqref{cylSchur}| as ``a cylindric or affine version of the Murnaghan--Nakayama rule'').
This is strong evidence, \emph{not} proof, that the two renderings number alike. Throughout I cite
by label, \verb|\eqref{cylSchur}|, and never write ``Eq.~(145)'' as though I had read that rendering.
\end{remark}

\section{Conventions, and the dictionary}\label{sec:dict}

\subsection{The AKO side}

Following \cite[l.\,313--330, 462--472]{AKO}: $\CC_{x,y}$ is $\mathbb{Z}^2$ with $(0,0)$ identified
with $(x,-y)$. In matrix coordinates $(r,c)$ (row index increasing downwards in the drawing) the
identification is
\begin{equation}\label{eq:shift}
  (r,c)\ \sim\ (r+y,\,c+x),
\end{equation}
so the cylinder has $y$ rows and $x$ columns per period; in particular its column set is
$\mathbb{Z}/x\mathbb{Z}$. A \emph{cylindric diagram} $D$ is the region between two boundary loops;
concretely I record it as row intervals $[\mathrm{lo}_r,\mathrm{hi}_r]$, $r\in\mathbb{Z}$, both
weakly increasing in $r$ and satisfying $\mathrm{lo}_{r+y}=\mathrm{lo}_r+x$,
$\mathrm{hi}_{r+y}=\mathrm{hi}_r+x$. Set $n:=x+y$.

A \emph{ribbon} is a connected cylindric diagram with no $2\times2$ block; a \emph{loop ribbon} is a
ribbon with exactly $n$ boxes. The height of a non-loop ribbon is (number of rows it meets) $-1$,
equivalently its number of vertical edges; all loop ribbons have height $y$.
Peeling loop ribbons off the outer boundary writes a \emph{stacked ribbon} as
$R=L_1\cup\dots\cup L_\ell\cup F$ with $F$ a loop ribbon (\emph{pure}) or a non-loop ribbon
(\emph{non-pure}).
A \emph{cylindric semistandard Young tableau} (CSSYT) of shape $D$ is a filling whose entries weakly
increase as the column index increases and strictly increase as the row index increases
(i.e.\ upwards in the drawing); $s_D=\sum_{T\in\CSSYT(D)}x^T$.

\begin{lemma}[well-definedness]\label{lem:welldef}
For all $364$ stacked ribbons on the $11$ cylinders $\CC_{x,y}$ with $x+y\le 6$ and $|R|\le 8$, every
peeling yields the same triple $(\ell,\ \text{pure?},\ \hgt(F)\bmod 2)$. Hence the weight assigned
below does not depend on the peeling.
\end{lemma}

\begin{lemma}[the width is the whole cylinder]\label{lem:width}
Every loop ribbon occupies all $x$ cylindrical columns. Consequently
$\wdt(R)=x$ for every pure stacked ribbon $R$, independently of $R$.
\end{lemma}

\begin{proof}
Read a loop ribbon as a cyclic lattice path: traversing one period it takes $x$ steps right and $y$
steps up/down, returning to its translate by \eqref{eq:shift}. It therefore meets a contiguous block
of $x+1$ columns $c_0,\dots,c_0+x$ of the universal cover, whose reductions modulo $x$ exhaust
$\mathbb{Z}/x\mathbb{Z}$ (with $c_0$ and $c_0+x$ the same cylindrical column). A pure stacked ribbon
contains a loop ribbon, so it too meets all $x$. $\square$
\end{proof}

\subsection{The KP side, and the translation}

\cite{KP} work with boxed partitions $\bl$ in the $k\times(n-k)$ box, cylindric loops
$\{(i,\bl_i)\}$ with $\bl_{i+k}=\bl_i-(n-k)$, and the shape $\bl/d/\bm$ whose row $i$ carries the
columns $\bm_i+1,\dots,\bl_{i-d}+d$. Their \emph{circular ribbon} has $n$ boxes and height $k$.

\begin{proposition}[dictionary]\label{prop:dict}
Put $x=n-k$ and $y=k$. Then $\bl/d/\bm$ is the cylindric diagram on $\CC_{x,y}$ given by
\begin{equation}\label{eq:build}
  \mathrm{lo}_r=\bm_{k-r}+1,\qquad \mathrm{hi}_r=\bl_{k-r-d}+d,\qquad r=0,\dots,k-1,
\end{equation}
extended by \eqref{eq:shift}. Under this identification KP's circular ribbons are exactly AKO's loop
ribbons, and KP's cylindric ribbons are exactly AKO's stacked ribbons. Moreover
$|\bl/d/\bm|=|\bl|-|\bm|+nd$, and transposing the diagram realises $\bl\mapsto\bl'$,
$\bm\mapsto\bm'$ up to a rotation of the cylinder and a $180^\circ$ rotation of the plane (both of
which preserve $s_D$).
\end{proposition}

\begin{proof}[Proof and calibration]
The row reversal $i=k-r$ is forced: $\bl$ is weakly decreasing in $i$ while $\mathrm{hi}$ must be
weakly increasing in $r$. Periodicity: $r\mapsto r+k$ sends $i\mapsto i-k$ and
$\bl_{i-k}=\bl_i+(n-k)=\bl_i+x$, matching \eqref{eq:shift}. The size identity is the computation
$\sum_{i=1}^{k}(\bl_{i-d}+d-\bm_i)=|\bl|+d(n-k)+kd-|\bm|$. A circular ribbon has $(n-k)+k=x+y=n$
boxes and $k=y$ vertical edges, i.e.\ is a loop ribbon. \emph{Calibration:} for $d=0$ the identity
$s_{\bl/0/\bm}=s_{\bl/\bm}$ (the ordinary skew Schur function), asserted by \cite{KP} just below
\verb|\eqref{cylSchur}|, holds for every case tested, and the whole of \eqref{eq:build} is
independently confirmed by the $d=0$ run of Theorem~\ref{thm:kp} below ($3731$ coefficients, $0$
mismatches). $\square$
\end{proof}

\begin{remark}[the dictionary is where the trap was]
It is tempting to read \eqref{cylSchur} as the assertion $\omega(s_D)=s_{D'}$, since \cite{KP} carry
both halves of $\omega$ explicitly (the factor $\epsilon_\nu$ and the conjugates $\bl',\bm'$).
That reading is \emph{wrong}, and it is refuted by the smallest cylindric shape there is:
on $\CC_{1,1}$ the two-box diagram $\mathrm{lo}=(0)$, $\mathrm{hi}=(1)$ is self-conjugate and
$s_D=e_2$, while $\omega(e_2)=h_2$. Testing $357$ shapes on nine cylinders, $\omega(s_D)=s_{D'}$
fails for $67$. It is exactly this failure that makes the two rules inequivalent (\S\ref{sec:verdict}).
\end{remark}

\section{The reduction: one number per region}

\begin{proposition}\label{prop:reduce}
Suppose a weight function $w$ on cylindric diagrams satisfies an MN rule
\begin{equation}\label{eq:mn}
  s_D=\sum_{\nu}\frac{p_\nu}{z_\nu}\sum_{\substack{\emptyset=D_0\subset D_1\subset\dots\subset D_{\ell}=D\\ |D_j/D_{j-1}|=\nu_j}}\ \prod_{j}w(D_j/D_{j-1})
\end{equation}
for all cylindric diagrams $D$. Then for every cylindric diagram $R$ with $|R|=r$,
\[
  w(R)=\langle s_R,\,p_r\rangle .
\]
In particular $w$ is unique, and any two candidate MN rules may be compared region by region.
\end{proposition}

\begin{proof}
Apply \eqref{eq:mn} to $D=R$ and $\nu=(r)$. A chain with a single step of size $r$ is
$\emptyset\subset R$ and there is exactly one, so the coefficient of $p_r/z_r$ in $s_R$ is $w(R)$.
Since $s_R=\sum_\nu c_\nu p_\nu/z_\nu$ and $\langle p_\nu,p_\rho\rangle=z_\nu\delta_{\nu\rho}$, that
coefficient is $\langle s_R,p_r\rangle$. $\square$
\end{proof}

This is the instrument. The right-hand side is computed from the CSSYT definition alone, so it is an
arbiter independent of both papers. As a first check it reproduces \cite[\texttt{ex:mnrule}]{AKO}
exactly (Example~\ref{ex:ako} below).

\section{Erratum A: the height of a stacked ribbon}

\cite[l.\,1182--1187]{AKO} give the height of a stacked ribbon twice:
\begin{quote}
``\dots by setting the height of a stacked ribbon $D=L_1\cup L_2\cup\dots\cup L_\ell\cup F$ to be
\emph{the sum of the heights of its constituents} (We can equivalently define the height as
\emph{the number of vertical edges of an extended ribbon} wrapping around the shape multiple times
and cylindrically covers the entire shape\dots) As all loop ribbons have height $y$, we have
$\hgt(D)=\ell y+\hgt(F)$.''
\end{quote}
The two definitions are \emph{not} equivalent. An extended ribbon is a path: on top of the vertical
edges internal to the constituents it must traverse one edge at each of the $\ell$ junctions, and
those edges are vertical. So the second definition exceeds the first by $\ell$, and the closed
formula quoted from the first is off by $\ell$.

\begin{theorem}[the weight function]\label{thm:weight}
For a cylindric diagram $R$ on $\CC_{x,y}$,
\[
  \langle s_R,p_{|R|}\rangle=
  \begin{cases}
    0, & R \text{ not a stacked ribbon},\\
    (-1)^{\,\ell(y+1)+\hgt(F)}\cdot 1, & R=L_1\cup\dots\cup L_\ell\cup F \text{ non-pure},\\
    (-1)^{\,\ell(y+1)+y}\cdot x, & R \text{ pure}.
  \end{cases}
\]
Equivalently: AKO's \texttt{prop:stackedRibbon} and \texttt{thm:cylindricMNRule} are correct verbatim
once $\hgt$ is read as the extended-ribbon count $\ell(y+1)+\hgt(F)$ rather than $\ell y+\hgt(F)$.
\end{theorem}

\emph{Status: verified exhaustively in range, and derived from \cite{K19} in \S\ref{sec:korff};
not proved here from first principles.} Over all $797$ regions on the $13$ cylinders $\CC_{x,y}$ with
$x+y\le6$ and $|R|\le8$ (of which $59$ are pure and $127$ have $\ell\ge1$): $0$ mismatches. The
negative controls all break, as they must:

\begin{center}
\begin{tabular}{lr}
\hline
weight variant & mismatches / $797$\\
\hline
Theorem~\ref{thm:weight} & \textbf{0}\\
$\hgt=\ell y+\hgt(F)$ (as stated in \cite{AKO}) & 174\\
Theorem~\ref{thm:weight} with $\wdt\equiv1$ (drop the width) & 44\\
Theorem~\ref{thm:weight} with the global sign flipped & 547\\
\hline
\end{tabular}
\end{center}

That the ``drop the width'' row is nonzero is the negative control certifying that the width factor is
load-bearing in the tested range: it is not an artefact of small $x$.

Feeding Theorem~\ref{thm:weight} into \eqref{eq:mn} reproduces the full multi-part rule:
over all cylindric diagrams with $x+y\le5$, $|D|\le6$, all $874$ pairs $(D,\nu)$ agree, $0$ mismatches.

\begin{example}[\cite{AKO}'s own example contradicts \cite{AKO}'s own formula]\label{ex:ako}
Let $D$ be \cite[\texttt{ex:mnrule}]{AKO}: on $\CC_{2,2}$, $\mathrm{lo}=(1,2)$, $\mathrm{hi}=(3,3)$,
so $|D|=5$ and $n=4$. Computing $s_D$ from the CSSYT definition and expanding:
\[
  s_D=-\frac{p_5}{z_5}+2\frac{p_{41}}{z_{41}}-2\frac{p_{311}}{z_{311}}+4\frac{p_{11111}}{z_{11111}},
\]
which is \emph{exactly} the expansion printed in \cite{AKO} --- so the example is right, and the
arbiter is calibrated. Now $D$ is a non-pure stacked ribbon with $\ell=1$ and $F$ a single box,
$\hgt(F)=0$. The stated formula gives $\hgt(D)=\ell y+\hgt(F)=2$, hence coefficient $(-1)^2=+1$ for
$p_5/z_5$. The truth is $-1$. The corrected formula gives $\ell(y+1)+\hgt(F)=3$, hence $-1$ ---
and $3$ is precisely the height \cite{AKO} annotate this very ribbon with (``Non-pure ribbon with
height 3''). The prose formula and the worked example disagree, and the example is the correct one.
\end{example}

\begin{remark}[why it survived]
The discrepancy is $(-1)^\ell$. It is invisible for $\ell=0$, i.e.\ for every ribbon that does not
wind --- which includes the entire classical specialisation, AKO's cancellation-free staircase
corollary as far as it is exercised by non-winding shapes, and four of the five stacked ribbons in
their illustrative example list. A parenthetical ``we can equivalently define\dots'' is graded by
nothing.
\end{remark}

\section{Erratum B: the missing multiplicity in \texorpdfstring{$\chi$}{chi}}

\cite[\texttt{\eqref{cylSchur}}]{KP} define
\begin{equation}\label{eq:kp}
  s_{\bl'/d/\bm'}=\sum_{\nu\in\mathcal{P}^+}\frac{\chi_{\bl/d/\bm}(\nu)}{z_\nu}\,\epsilon_\nu\,p_\nu,
  \qquad \epsilon_\nu=(-1)^{|\nu|-\ell(\nu)},
\end{equation}
with $\chi_{\bl/d/\bm}(\nu)=\sum_{\bar\pi}\prod_{i\ge0}(-1)^{\hgt(\bar\pi^{-1}(i))-1}$ over cylindric
ribbon plane partitions of shape $\bl/d/\bm$ and weight $\nu$, and prove that \eqref{eq:kp} equals the
cylindric Schur function of the conjugate shape. Their height is the number of rows occupied, with
consecutive pieces sharing a row, so for a cylindric ribbon with $\ell$ windings and tail $F$
\begin{equation}\label{eq:kpht}
  \hgt^{\mathrm{KP}}=\bigl(\hgt(F)+1\bigr)+\ell(k-1),\qquad
  \hgt^{\mathrm{KP}}=m(k-1)+1 \text{ for } m \text{ circular ribbons and no tail,}
\end{equation}
which is exactly their recurrence $\chi_{\bl/d/\bm}(r)=(-1)^{k-1}\chi_{\bl/(d-1)/\bm}(r-n)$ and their
Figure caption ($\hgt=2+2-1=3$ for a ribbon of length $5+4$).

\begin{lemma}[the two signs already agree]\label{lem:signs}
With $y=k$, for every stacked ribbon $R$,
\[
  (-1)^{\hgt^{\mathrm{KP}}(R)-1}=
  \begin{cases}
    (-1)^{\ell(y+1)+\hgt(F)}, & R \text{ non-pure},\\
    -\,(-1)^{\ell(y+1)+y}, & R \text{ pure}.
  \end{cases}
\]
\end{lemma}

\begin{proof}
Non-pure: $\hgt^{\mathrm{KP}}-1=\hgt(F)+\ell(y-1)$ and
$\ell(y+1)-\ell(y-1)=2\ell$ is even. Pure with $m=\ell+1$ loops:
$\hgt^{\mathrm{KP}}-1=m(y-1)$ while $\ell(y+1)+y=m(y+1)-1$, and
$m(y+1)-1-m(y-1)=2m-1$ is odd. $\square$
\end{proof}

So, once Erratum A is applied, \emph{the sign conventions of \cite{AKO} and \cite{KP} reconcile
perfectly on non-winding-tail ribbons and differ by exactly one sign on fully wound ones.} All that
is left between them is a multiplicity.

\begin{theorem}[what \eqref{eq:kp} needs]\label{thm:kp}
Let $D=\bl/d/\bm$ on $\CC_{n-k,k}$ and $E=\bl'/d/\bm'$ on $\CC_{k,n-k}$. Define
$\tilde\chi_D(\nu)$ by \eqref{eq:mn}-style chains on $D$ with per-step weight
\[
  v(R)=(-1)^{\hgt^{\mathrm{KP}}(R)-1}\cdot
  \begin{cases} k, & R \text{ fully wound (all circular ribbons)},\\ 1,&\text{otherwise.}\end{cases}
\]
Then $\tilde\chi_D(\nu)=\epsilon_\nu\,\langle s_E,p_\nu\rangle$, i.e.\ \eqref{eq:kp} holds with
$\chi$ replaced by $\tilde\chi$.
\end{theorem}

\emph{Status: verified, not proved.} Over all $(\bl,\bm,d)$ with $n\le7$, $0\le d\le3$ and
$|D|\le7$ on $14$ pairs $(n,k)$:

\begin{center}
\begin{tabular}{lrr}
\hline
multiplicity on a fully wound step & coefficients tested & mismatches\\
\hline
$k$ (Theorem~\ref{thm:kp}) & $16\,791$ & \textbf{0}\\
$1$ (\cite{KP} as literally defined) & $16\,791$ & 439\\
$n-k$ & $16\,791$ & 425\\
\hline
$1$, restricted to $d=0$ (calibration) & $3\,731$ & \textbf{0}\\
\hline
\end{tabular}
\end{center}

The last row is the calibration demanded of any such comparison: at $d=0$ nothing winds, the
multiplicity is invisible, and \eqref{eq:kp} reduces to the classical
$s_{\lambda'/\mu'}=\sum_\nu\epsilon_\nu\chi_{\lambda/\mu}(\nu)p_\nu/z_\nu$ --- which passes. So the
dictionary of Proposition~\ref{prop:dict} is right and the $d\ge1$ failures are real.

\begin{erratum}\label{err:kp}
$\chi_{\bl/d/\bm}(\nu)$ as defined in \cite{KP} is short by a factor $k$ for each part of $\nu$
consumed by a fully wound cylindric ribbon. Minimal instance: $n=4$, $k=2$, $\bl=\bm=(0,0)$, $d=1$,
so $D$ is a \emph{single circular ribbon} of length $4$ on $\CC_{2,2}$, and $\nu=(4)$. The
definition gives one ribbon plane partition of height $\hgt^{\mathrm{KP}}=k=2$, hence
$\chi=(-1)^{2-1}=-1$. But $E=D$ here and $\langle s_E,p_4\rangle=+2$ with
$\epsilon_{(4)}=-1$, so \eqref{eq:kp} requires $\chi=-2$.
\end{erratum}

\begin{remark}[where the factor went]
The multiplicity is present in \cite{KP}'s \emph{proof} and absent only from the combinatorial
definition. Their Lemma is stated for the matrix element
$\langle v^{\bl},P^*_\nu v_{\bm}\rangle$, and the proof reads ``$P^*_r$ adds \emph{all possible}
ribbons of length $r$ to $\bm$''. Their \texttt{lem:CR} asserts a \emph{unique} starting diagonal
$a$ only under the hypothesis $r\notin n\mathbb{N}$; when $r\in n\mathbb{N}$ the starting diagonal is
not unique, several $a$ give the same shape, and the matrix element counts them. A cylindric ribbon
\emph{plane partition}, however, records only shapes --- the level $\bar\pi^{-1}(i)$ is a region, not
a region-with-a-marked-starting-diagonal --- so the degeneracy is lost in passing from the matrix
element to $\chi$. The hypothesis $r\notin n\mathbb{N}$ in \texttt{lem:CR} is where the information
was, and it is exactly the case excluded there that carries the multiplicity.
\end{remark}

\section{Reconciliation with Korff (2019), and a proof of Erratum A}\label{sec:korff}

Korff \cite[\texttt{lem:cylMNrule}]{K19} proves a cylindric MN rule for $t$-deformed cylindric Hecke
characters $\chi^{\lambda/d/\mu}_t$, defined by matrix elements at the \emph{conjugates}
$\lambda',\mu'$ and satisfying $\mathrm{ch}_t(\chi^{\lambda/d/\emptyset}_t)=s_{\lambda/d/\emptyset}$
(\cite[\texttt{cor:cylchi2cylschur}]{K19}) --- so there the ribbon combinatorics and the Schur
function live on the \emph{same} cylinder $\CC_{n-k,k}$. In the $t\to1$ limit
\cite[Remark following \texttt{lem:cylMNrule}]{K19} his two clauses read
\begin{align}
  \text{(i)}\quad & \chi^{\lambda/d/\mu}(\nu_1,\dots,\nu_i+n,\dots,\nu_\ell)=(-1)^{k-1}\,\chi^{\lambda/d-1/\mu}(\nu),
  \label{eq:korffi}\\
  \text{(ii)}\quad & \chi^{\lambda/d/\mu}(\nu,n)=(-1)^{k}(n-k)\,\chi^{\lambda/d-1/\mu}(\nu).
  \label{eq:korffii}
\end{align}

\begin{corollary}\label{cor:korff}
Theorem~\ref{thm:weight} is exactly equivalent to \eqref{eq:korffi}--\eqref{eq:korffii}, and
\cite{AKO}'s stated formula $\hgt=\ell y+\hgt(F)$ is not.
\end{corollary}

\begin{proof}
With $y=k$ and $x=n-k$, Theorem~\ref{thm:weight} gives the weight
$(-1)^{\ell(y+1)+\hgt(F)}$, so raising $\ell$ by one --- adding one winding, which is
\eqref{eq:korffi} --- multiplies by $(-1)^{y+1}=(-1)^{k-1}$: clause (i). The $\ell=0$ pure case, a
single loop ribbon, has weight $(-1)^{y}x=(-1)^k(n-k)$: clause (ii). Under the stated formula
$\hgt=\ell y+\hgt(F)$ the first factor would instead be $(-1)^{y}=(-1)^{k}$, contradicting
\eqref{eq:korffi}. $\square$
\end{proof}

Since \eqref{eq:korffi} is proved in \cite{K19} directly from the operator identity
$H_{\nu_i+n}=(-1)^{k-1}q\,H_{\nu_i}$, Corollary~\ref{cor:korff} upgrades Erratum A from
``verified on $797$ regions'' to a consequence of a proved lemma in the literature.

\begin{remark}[novelty gate]
The corrected AKO rule is therefore, in the $t\to1$ limit, \cite[\texttt{lem:cylMNrule}]{K19}
(2019) --- which predates \cite{AKO} and which \cite{AKO} do not cite at this point; they cite
\cite{KP} instead, and \cite{KP} is the \emph{further} of the two relatives. What is new here is
neither rule: it is (a) the dictionary of Proposition~\ref{prop:dict} with its $d=0$ calibration,
(b) the reduction of Proposition~\ref{prop:reduce} that makes the comparison decidable, (c) the two
errata with minimal counterexamples, and (d) the verdict below. \cite{AKO}'s own contribution
--- a combinatorial proof and a closed, non-recursive form --- is untouched by Erratum A, which is a
defect in one formula, not in the theorem.
\end{remark}

\section{Verdict on the hedge}\label{sec:verdict}

\begin{theorem}[the two rules are different theorems]\label{thm:verdict}
Write $\chi^{A}_D(\nu)$ for the summand total of \cite[\texttt{thm:cylindricMNRule}]{AKO} on a
cylindric diagram $D$ (with the corrected height), and $\tilde\chi^{K}_D(\nu)$ for that of
Theorem~\ref{thm:kp}. Then
\[
  \chi^{A}_D(\nu)=\langle s_D,p_\nu\rangle,\qquad
  \tilde\chi^{K}_D(\nu)=\epsilon_\nu\,\langle s_{D'},p_\nu\rangle,
\]
so $\chi^{A}=\tilde\chi^{K}$ would force $\omega(s_D)=s_{D'}$ for all cylindric $D$ --- which is
false ($67$ failures in $357$ shapes; the minimal one is $s_D=e_2$ on $\CC_{1,1}$, self-conjugate).
Concretely the two rules sum over the \emph{same} sets of stacked $=$ cylindric ribbons with the
\emph{same} signs (Lemma~\ref{lem:signs}), and differ precisely on fully wound steps, by the factor
\[
  \frac{v(R)}{w(R)}=-\frac{k}{\,n-k\,}.
\]
\end{theorem}

The hedge ``a similar formula seems to appear'' is thus \emph{correct as a hedge and false as an
identification}. The two statements are the two sides of level-rank duality
$qH^*(\mathrm{Gr}(k,n))\cong qH^*(\mathrm{Gr}(n-k,n))$: the multiplicity on a fully wound ribbon is
the width of the cylinder on which the \emph{Schur function} lives, $n-k$ for \cite{AKO} and
\cite{K19} (ribbons and Schur function on one cylinder), $k$ for \cite{KP} (ribbons on one cylinder,
Schur function on the conjugate). Classically the two conventions are interchangeable because
$\omega(s_{\lambda/\mu})=s_{\lambda'/\mu'}$; cylindrically $\omega$ no longer conjugates the diagram,
and the two rules part company exactly where the ribbon winds.

\section{Gaps}

\begin{enumerate}
\item \textbf{Theorem~\ref{thm:weight} and Theorem~\ref{thm:kp} are verified, not proved from first
principles.} Ranges: $797$ regions / $874$ multi-part pairs / $16\,791$ KP coefficients, all with
$0$ mismatches, on $x+y\le6$ resp.\ $n\le7$, $|D|\le8$, $d\le3$. Erratum A is additionally implied by
\cite[\eqref{eq:korffi}]{K19} (Corollary~\ref{cor:korff}); Erratum B is not --- it rests on the
computation plus the structural reason given in the Remark after Erratum~\ref{err:kp}.
\item \textbf{The extended-ribbon reading of Erratum A is an explanation, not a proof.} I verified
that AKO's two definitions differ by $\ell$ as a closed formula, and that each junction between
consecutive constituents carries at least one vertical edge; but I did \emph{not} prove that an
extended ribbon always exists, nor that it must turn vertically at every junction. In $37$ of $127$
tested stacked ribbons with $\ell\ge1$ the junction carries \emph{two} vertical adjacencies, of which
a path uses one --- consistent with, but not a proof of, the count $\ell(y+1)+\hgt(F)$.
\item \textbf{Equation numbering.} ``Eq.~(145)'' is identified with \verb|\eqref{cylSchur}| by an
environment count in the e-print, not by reading the published rendering (see Remark~1.1).
\item \textbf{Corollary~\ref{cor:korff} leans on an implicit extension.}
\cite[\texttt{cor:cylchi2cylschur}]{K19} is stated for the \emph{non-skew} case
$\mathrm{ch}_t(\chi^{\lambda/d/\emptyset}_t)=s_{\lambda/d/\emptyset}$. I have used it in the skew
form $\chi^{\lambda/d/\mu}(\nu)=\langle s_{\lambda/d/\mu},p_\nu\rangle$, which \cite{K19} treats as
understood but does not display. If that extension fails, Erratum~A reverts from ``proved'' to
``verified on $797$ regions'', which is still its status independently.
\item \textbf{Lemma~\ref{lem:welldef}} (uniqueness of the peeling invariants) is verified in range,
not proved.
\end{enumerate}

\begin{thebibliography}{9}
\bibitem{AKO} P. Alexandersson, E. Kantarc\i\ O\u{g}uz, \emph{Cylindric Schur functions},
\texttt{arXiv:2311.07382}; source file \texttt{Cylindric-Schur-functions-July-merge.tex}, cited by
internal label and source line (\texttt{thm:cylindricMNRule} l.\,1269, \texttt{prop:stackedRibbon}
l.\,1240, \texttt{ex:mnrule} l.\,1284, height discussion l.\,1182, hedge l.\,232,
\texttt{def:cylindricSkewSchur} l.\,752).
\bibitem{KP} C. Korff, D. Palazzo, \emph{Cylindric Hecke characters and Gromov--Witten invariants
via the asymmetric six-vertex model}, \texttt{arXiv:1804.05647} (published as Algebr. Comb.
\textbf{3} (2020) 191--247); cited by internal label (\verb|\eqref{cylSchur}| l.\,1955, the matrix
element Lemma l.\,1928 and its proof l.\,1940--1955, \texttt{lem:CR} l.\,1922,
$\epsilon_\nu$ l.\,489).
\bibitem{K19} C. Korff, \emph{(cylindric Hecke characters / asymmetric six-vertex model)},
\texttt{arXiv:1906.02565}, draft \texttt{korff-1906.02565-draft\_02.tex}; \texttt{lem:cylMNrule}
l.\,1770, its $t\to1$ Remark l.\,1831, \texttt{cor:cylchi2cylschur} l.\,2037. Cited by internal
label and source line, per the same convention.
\end{thebibliography}

\end{document}
